\begin{table}[htbp]
\centering
\caption{Sensitivity analyses for the nine genes carrying cis-pQTL instruments.}
\label{tab:mr_sensitivity}
\footnotesize
\begin{tabular}{lllll}
\toprule
\multicolumn{5}{l}{\textbf{(A)} MR method comparison}\\
\midrule
  Gene & $n$ & IVW $b$ ($p$) & Weighted median $b$ ($p$) & Mode $b$ ($p$) \\
\midrule
  TNFRSF1A  & 25    & $-0.013$ (0.637)                & $+0.038$ (0.034)                 & $+0.030$ (0.075) \\
  NCF2      & 28    & $-0.033$ (0.156)                & $-0.105$ ($6.57\times10^{-6}$)   & $-0.115$ ($1.06\times10^{-4}$) \\
  STAT6     & 61    & $-0.007$ (0.427)                & $-0.012$ (0.082)                 & $-0.017$ ($7.37\times10^{-4}$) \\
  LIMA1     & 70    & $+0.882$ ($1.34\times10^{-26}$) & $+0.966$ ($7.03\times10^{-104}$) & $+0.909$ ($7.32\times10^{-66}$) \\
  TNFRSF1B  & 73    & $-0.051$ ($3.03\times10^{-7}$)  & $-0.042$ ($1.94\times10^{-3}$)   & $-0.040$ (0.096) \\
  CCM2      & 94    & $-0.515$ ($2.72\times10^{-17}$) & $-0.719$ ($7.26\times10^{-133}$) & $-0.734$ ($2.20\times10^{-44}$) \\
  CDKN1A    & 228   & $-0.131$ ($1.28\times10^{-18}$) & $-0.111$ ($1.86\times10^{-19}$)  & $-0.116$ ($1.41\times10^{-8}$) \\
  TNF$^{*}$ & 2288  & $-0.062$ ($8.14\times10^{-29}$) & $-0.014$ (0.085)                 & $+0.001$ (0.867) \\
  LTA$^{*}$ & 10213 & $-0.022$ ($5.41\times10^{-44}$) & $-0.012$ ($3.39\times10^{-22}$)  & $-0.001$ (0.907) \\
\bottomrule
\end{tabular}

\vspace{1.4em}

\begin{tabular}{lllll}
\toprule
\multicolumn{5}{l}{\textbf{(B)} Pleiotropy and heterogeneity diagnostics}\\
\midrule
  Gene & Egger intercept ($p$) & Cochran $Q$ $p$ & MR-PRESSO $p$ & MR-PRESSO distortion $p$ \\
\midrule
  TNFRSF1A  & $-0.0169$ ($5.13\times10^{-4}$)  & $2.28\times10^{-10}$  & $<2.0\times10^{-4}$ & 0.505 \\
  NCF2      & $+0.0409$ (0.034)                & $7.86\times10^{-11}$  & $2.0\times10^{-4}$ & $<2.0\times10^{-4}$ \\
  STAT6     & $+0.0205$ ($3.04\times10^{-12}$) & $1.74\times10^{-13}$  & $<2.0\times10^{-4}$ & 0.085 \\
  LIMA1     & $+0.0744$ (0.109)                & $5.28\times10^{-81}$  & $<2.0\times10^{-4}$ & $<2.0\times10^{-4}$ \\
  TNFRSF1B  & $-0.0138$ ($8.67\times10^{-5}$)  & 0.997                 & 0.997 & ---$^{\S}$ \\
  CCM2      & $+0.0129$ (0.549)                & $2.76\times10^{-107}$ & $<2.0\times10^{-4}$ & $<2.0\times10^{-4}$ \\
  CDKN1A    & $-0.0122$ ($1.76\times10^{-4}$)  & $8.78\times10^{-48}$  & $<2.0\times10^{-4}$ & $<2.0\times10^{-4}$ \\
  TNF$^{*}$ & $-0.0117$ ($<10^{-300}$)         & $<10^{-300}$          & skipped$^{\ddagger}$ & skipped$^{\ddagger}$ \\
  LTA$^{*}$ & $-0.0018$ ($7.41\times10^{-4}$)  & $<10^{-300}$          & skipped$^{\ddagger}$ & skipped$^{\ddagger}$ \\
\bottomrule
\end{tabular}

\vspace{1.4em}

\begin{tabular}{lll}
\toprule
\multicolumn{3}{l}{\textbf{(C)} Outlier handling}\\
\midrule
  Gene & Outliers & Outlier-corrected $b$ ($p$) \\
\midrule
  TNFRSF1A  & 5 & $-0.086$ (0.028) \\
  NCF2      & 3 & $+0.047$ (0.004) \\
  STAT6     & 5  & $+0.004$ (0.666) \\
  LIMA1     & 10 & $+1.076$ ($8.06\times10^{-37}$) \\
  TNFRSF1B  & none$^{\S}$ & ---$^{\S}$ \\
  CCM2      & 15 & $-0.762$ ($1.90\times10^{-57}$) \\
  CDKN1A    & 22 & $-0.102$ ($5.91\times10^{-20}$) \\
  TNF$^{*}$ & --- & skipped$^{\ddagger}$ \\
  LTA$^{*}$ & --- & skipped$^{\ddagger}$ \\
\bottomrule
\end{tabular}

\tblnote{\textbf{(A)} Effect estimates across MR methods, on the log-odds scale per 1 SD of genetically predicted protein level; genes are listed in ascending order of instrument count. \textbf{(B)} Pleiotropy and heterogeneity diagnostics: Cochran's $Q$ $p$-values test for heterogeneity across instruments; a non-significant MR-Egger intercept indicates no detectable directional pleiotropy; MR-PRESSO was run with \texttt{NbDistribution = 5000} so that the outlier test resolves $n/\mathrm{NbDistribution} \le 0.046$ for every gene tested, whereas at the package default of 1000 the same quantity is 0.061--0.228, above the 0.05 significance threshold, so the default cannot identify outliers at this instrument count. \textbf{(C)} Outlier handling: the outlier-corrected column reports the IVW estimate after removing the instruments flagged by the MR-PRESSO outlier test, and the Outliers column gives their number. Outlier-test $p$-values that fall at the resolution floor ($n/\mathrm{NbDistribution}$) are reported below as inequalities and are counted as outliers. For NCF2, removing 3 of the 28 instruments reverses the sign of the estimate (Table 3A versus 3C), so the direction of effect is not interpreted for that gene. $^{\S}$MR-PRESSO's outlier and distortion tests are performed only when the global test is significant, so for TNFRSF1B (p = 0.997) no instrument was tested for outlyingness and no outlier-corrected estimate exists. $^{*}$Located in the MHC region (chr6:29--33 Mb), where strong long-range LD violates the instrument-independence assumption of every method shown; estimates are reported for completeness and should be read as reference values only. $^{\ddagger}$Not computationally feasible: MR-PRESSO requires more than $n/0.046$ draws (about $5.0\times10^{4}$ for $n = 2288$ and $2.2\times10^{5}$ for $n = 10213$).}
\end{table}